\subsubsection{Spectrum in the presence of brane-localized mass on torus extra dimensions --- arXiv:1607.07152}

\textbf{Authors:} Yutaka Sakamura

\paragraph{主な主張}
6次元理論を $T^2$ にコンパクト化してbrane-localized massを導入すると、最軽量質量固有値はcutoff $\Lambda$ に対して無視できない依存性を持ち、$\Lambda$ がcompactification scaleより2桁高くてもthin-brane近似のUV感度が残る一方、その固有値はtorus modulus $\tau$ にも依存し、大きなbrane mass極限でも典型的にcompactification scaleより1桁程度軽いことを示した\cite{Sakamura:2016BraneTorus}。

\paragraph{新規性}
5次元では弱いbrane-localized massのUV感度がcodimension 2の$T^2$では顕著になることを数値対角化と近似式で定量化し、cutoff依存とcomplex-structure modulus $\tau$ 依存を同時に明示した。

\paragraph{位置付け}
$T^2$ bulk-neutrino模型で現れる対数的UV感度を理解する直接的な基礎であり、2次元KK格子では $N_{\rm KK}(\Lambda)\propto \Lambda^2$ なので $\sum_{\rm KK}1/M_{\rm KK}^2\sim\ln\Lambda$ となることから、pointlike brane近似の有限厚みやUV completionを検討する動機を与える。
